\section{Глава 1. Одно-агентная задача}

В рамках первого эксперимеента рассмотрим сценарий задачи с одним агентом - такой формат обучения 
будет освобожден от ограничений взаимодействия с другими агентами в виду их отсутствия, 
и позволит максимально пристально оценить влиение шока и эффективность послешокового обучения.

\subsection{Сценарий}

Эксперименты проводятся для одного самолета, движущегося из каждого входа слева со следующим сценарием.

\begin{enumerate}
    \item Сначала базовый агент обучается движению от ворот до ближайшего к воротам выходу - 
		по одному отдельному агенту на ворота 
    \item После обучения, проводится эксперимент окружения на целевой сети с введением шока:
		на пути к предобученному пункту назначения, назначенный выход закрывается и проводится 
		попытка продолжить выполнение эпизода - новым пунктом назначения выбирается последний 
		оставшийся выход.
		Шаг введение шока в окружение можно варьировать, но для экспериментов выбираются такие моменты,
		чтобы после введения шока у агента оставалось пространство для маневрирования.
    \item Инициализуруется процесс дообучения модели с тёплым стартом (с параметрами базовой модели 
		в качестве начальных) с задачей перемещения самолёта от точки, в которой произошёл шок, до нового пункта назначения.
\end{enumerate}

\subsection{Архитектура сети и параметры}

Состояние окружающей среды декодируется как изображение с двумя слоями 
– картой и положением плоскости (см. рисунок \ref{fig:state_map}), 
итоговая форма изображения - $2 \times 9 \times 10$.

Нейронная сеть состоит из следующих слоев:

\begin{itemize}
    \item Сверточный слой с фильтром размера 2x3x3, одним входным и 8 выходными слоями 
    \item Сверточный слой с фильтром размера 3x3,  8 входными и 16 выходными слоями 
    \item Сверточный слой с фильтром размера 3x3, 16 входными и 32 выходными слоями 
    \item Сверточный слой с фильтром размера 3x3, 32 входными и 64 выходными слоями 
    \item Полносвязный слой с 128 входами and 6 выходами (равно количеству доступных действий)
\end{itemize}


\begin{table}[h]
    \centering
    \begin{tabular}{|c|c|} 
        \hline
        Объём обучаещего сета $N_B$ & 500 \\ \hline
        дисконт. коэф-ент $\gamma$ & 0.99 \\ \hline
        функция потерь $L_Q$ & Huber loss \\ \hline
        алгоритм оптимизации & Adam \\ \hline
        коэф-ент оптимизации $\beta$ & $10^{-4}$ \\ \hline
        $n_{r}$ & 500 \\
        \hline
    \end{tabular}
    \label{tab:fix_hyp}
    \caption{Параметры обучения моделей}
\end{table}

\subsection{Результаты}


Результаты эксперимента продемонстрированы на рисунке \ref{fig:retrain_res} в приложении.
В каждой строке показаны результаты обучения и переобучения одного агента из разных ворот.

Траектории базовых обученных агентов изображены в левом столбце.
Каждый агент выполнил свои задачи и сумел добраться до места назначения (самолёты из двух верхних ворот добрались до верхнего выходу, 
а из двух нижних ворот - к нижнемму выходу).

В средней колонке траектория выполнялась по нетронутой карте окружения до красного креста, где находился самолет, 
когда произошел шок и закрылись ворота, запланированные как пункт назначения.
Последний доступный выход становится новым пунктом назначения, и эксперимент продолжается.
Следующие шаги агента приводили к сбою (первый и третий самолеты наступили на неразрешенную ячейку) или 
к завершению эпизода по истечению времени (второй самолет остановился, а четвертый самолет шел по карте, пока не достиг предела шага).

Траектории переобученных моделей показаны в правых столбцах.
Задача переобученного агента заключалась в том, чтобы проследовать к последнему доступному выходу,
начиная из точки, где произошел шок (красный крест), а не со стартовых ворот.
Каждому агенту удалось добраться до другого выхода и выполнить задачу после того, как произошел шок окружения.

Из таблицы \ref{tab:retrain_vs} хорошо видно, что процесс послешокового дообучения занял значительно меньше эпизодов и воспроизводящей памяти,
чем первоначальное обучение, что видно из рисунков \ref{fig:rew_graph12} и \ref{fig:rew_graph34}.

Это можно объяснить отсутствием принципиальных различий (как например в системе наград или пространстве действий агента) между базовой и
измененной окружающей средой и сохранением фундаментальных «знаний» о правилах окружающей среды в процессе дообучения.

\begin{table}[h]
    \centering
    \begin{tabular}{|c|c|c|}
        \hline
        этап обучения & базовый & послешоковый \\ \hline
        число эпизодов & 5000 & 1500 \\ \hline
        объём ВП & 10000 & 1000 \\ \hline
        ур. обучения целовой сети $\tau$ & 0.1 & 0.2 \\ \hline
    \end{tabular}
    \caption{Гиперпараметры обучения и базовой и послешоковой модели}
    \label{tab:retrain_vs}
\end{table}

Результаты эксперимента показывают временную (при значительно меньшем количестве эпизодов) и ресурсную
эффективность (за счет меньшего объема памяти повторов) метода постшокового дообучения даже при
минимальной настройке гиперпараметров обучения агентов в случае экологических синтетических шоков.